\section{Maquettes et enchaînement des écrans}
\label{sec:maquettes}

L'application comporte neuf écrans : connexion, inscription, fil d'actualité, détail d'un
message, profil utilisateur, édition de profil, recherche, administration et page d'erreur 404.
Leur enchaînement est formalisé par la figure~\ref{fig:enchainement}.

\begin{figure}[ht]
\centering
\resizebox{\textwidth}{!}{% Enchaînement des écrans (storyboard / plan de navigation)
\begin{tikzpicture}[
  node distance=10mm and 14mm,
  screen/.style={draw=solageBlue!70, fill=infoBg, rounded corners=2pt,
    align=center, font=\scriptsize, minimum width=2.4cm, minimum height=0.9cm},
  guard/.style={draw=juryGreen!70, fill=juryBg, rounded corners=2pt,
    align=center, font=\scriptsize, minimum width=2.4cm, minimum height=0.9cm},
  admin/.style={draw=todoOrange!70, fill=todoBg, rounded corners=2pt,
    align=center, font=\scriptsize, minimum width=2.4cm, minimum height=0.9cm},
  flow/.style={-{Stealth[length=2mm]}, semithick, solageGray}
]
\node[screen] (login)    {/login\\{\scriptsize Connexion}};
\node[screen, below=of login] (register) {/register\\{\scriptsize Inscription}};
\node[guard, right=24mm of login] (feed) {/\\{\scriptsize Fil d'actualité}};

\node[screen, right=of feed] (post) {/post/\{id\}\\{\scriptsize Détail + réponses}};
\node[screen, above=of post] (search) {/search\\{\scriptsize Recherche}};
\node[screen, below=of feed] (user) {/user/\{id\}\\{\scriptsize Profil}};
\node[screen, below=of user] (edit) {/edituser/\{id\}\\{\scriptsize Édition profil}};
\node[admin, below=of post] (admin) {/admin\\{\scriptsize Modération}};
\node[screen, draw=solageGray, fill=gray!10, right=of search] (e404) {404\\{\scriptsize Page inconnue}};

\draw[flow] (login) -- (register) node[midway,left,font=\scriptsize]{lien};
\draw[flow] (login) -- (feed) node[midway,above,font=\scriptsize]{auth OK};
\draw[flow] (register.east) to[out=0,in=-135] (feed);
\draw[flow] (feed) -- (post);
\draw[flow] (feed) -- (search);
\draw[flow] (feed) -- (user);
\draw[flow] (user) -- (edit);
\draw[flow] (feed) -- (admin) node[midway,right,font=\scriptsize]{si Admin};
\draw[flow] (post.north) -- (search.south);
\draw[flow, dashed] (feed) to[out=45,in=180] (e404);
\end{tikzpicture}
}
\caption{Enchaînement des écrans (storyboard). En vert l'écran protégé par authentification,
en orange l'écran réservé à l'administrateur.}
\label{fig:enchainement}
\end{figure}

J'ai délibérément réutilisé les \textbf{conventions d'interface éprouvées} du microblogging — fil
chronologique, zone de composition en haut, réponses imbriquées, action « j'aime » sous chaque
message — plutôt que d'inventer une interface nouvelle. Ces motifs étant déjà familiers aux
utilisateurs, les reprendre maximise l'\textbf{apprenabilité} et sert les principes ergonomiques
de simplicité et de minimalité des affichages. Ma contribution de conception tient dans
l'\textbf{adaptation} : un périmètre réduit et cohérent, une charte graphique minimaliste propre
à Solage, et l'enchaînement des écrans présenté ci-dessus. Le parcours nominal mène de la
connexion au fil d'actualité, depuis lequel l'utilisateur accède au détail d'un message, aux
profils, à la recherche et — s'il est administrateur — à la modération.

J'ai formalisé ces écrans sous forme de \textbf{maquettes}, réalisées avec \textbf{Penpot} (outil
de maquettage \emph{open source}). Ce sont mes propres écrans — identité visuelle, charte et
libellés propres à \projet{},
sans aucun élément de marque tiers. Chaque écran est décliné en version \textbf{bureau} et
\textbf{mobile}, actant une conception \emph{responsive}. Les trois écrans les plus structurants
sont présentés ici ; les autres figurent en annexe~\ref{ann:maquettes}.

\begin{figure}[H]
\centering
\includegraphics[width=0.82\textwidth]{desktop-connexion.png}
\caption{Maquette de l'écran de connexion (bureau).}
\label{fig:maq-connexion}
\end{figure}

\paragraph{Connexion — focaliser sur une tâche unique.} L'écran de connexion isole sa seule
tâche dans une \textbf{carte centrée} : logo et accroche au-dessus, deux champs étiquetés
explicitement (« Email », « Mot de passe »), et un \textbf{unique appel à l'action} mis en avant
par la couleur d'accent violette. L'absence de tout élément concurrent réduit la charge cognitive
et guide l'œil vers l'action attendue.

\begin{figure}[H]
\centering
\begin{minipage}[t]{0.45\textwidth}\centering
  
\includegraphics[height=8.6cm]{mobile-accueil.png}
\end{minipage}\hfill
\begin{minipage}[t]{0.45\textwidth}\centering
  
\includegraphics[height=8.6cm]{mobile-detail-post.png}
\end{minipage}
\caption{Maquettes mobiles : le fil d'actualité (à gauche) et le détail d'un message avec sa zone
de réponse (à droite).}
\label{fig:maq-mobile}
\end{figure}

\paragraph{Fil et détail — hiérarchie visuelle et constance.} Sur le fil comme sur le détail, la
\textbf{zone de composition} reste en haut, immédiatement accessible. Chaque message suit la même
\textbf{hiérarchie visuelle} : nom de l'auteur en gras, horodatage atténué en gris, contenu en
corps de texte, puis les actions (« répondre », « j'aime ») alignées sous le message. Sur mobile,
une \textbf{barre de navigation inférieure} (accueil, recherche, profil) reste constante d'un
écran à l'autre, et un bouton de retour ramène du détail vers le fil — deux repères de navigation
qui ne bougent jamais.

\paragraph{Choix d'accessibilité (RGAA).} Les maquettes intègrent les repères d'accessibilité
visés par le projet : libellés de formulaire explicites associés à chaque champ, contraste
suffisant entre le texte et le fond, et zones d'action (boutons, icônes) dimensionnées pour le
tactile. Ces choix de conception se traduisent dans le code par des attributs \texttt{alt},
\texttt{aria-label} et des libellés \texttt{<label>} (cf. chapitre~\ref{ch:realisations}).

\section{Conception de la base de données}
\label{sec:bdd}

\subsection{Modèle conceptuel des données (MCD)}

Le modèle conceptuel (figure~\ref{fig:mcd}) recense trois entités principales :
\textsc{Utilisateur}, \textsc{Role} et \textsc{Post}. Les réponses sont modélisées par une
association \emph{réflexive} « répond à » sur \textsc{Post} : une réponse \emph{est} un message
qui en référence un autre. Les « j'aime » et les favoris sont des associations entre
\textsc{Utilisateur} et \textsc{Post}, l'association « aime » portant l'attribut
\texttt{created\_at}.

\begin{figure}[ht]
\centering
\resizebox{0.95\textwidth}{!}{% MCD (Merise) — modèle conceptuel des données de Solage
\begin{tikzpicture}[
  entity/.style={draw=solageBlue!75, fill=infoBg, rounded corners=1pt,
    align=left, font=\small, inner sep=5pt},
  assoc/.style={draw=juryGreen!70, fill=juryBg, ellipse,
    align=center, font=\footnotesize\itshape, inner sep=1pt, minimum width=2cm, minimum height=1cm},
  card/.style={font=\scriptsize, text=solageDark, fill=white, inner sep=1pt},
  link/.style={solageGray, semithick}
]
% --- Entités ---
\node[entity] (role) at (0,0)
  {\textbf{ROLE}\\[2pt]\underline{id}\\ name};
\node[entity] (user) at (5.2,-0.3)
  {\textbf{UTILISATEUR}\\[2pt]\underline{id}\\ name\\ firstname\\ email\\ password\\ image};
\node[entity] (post) at (11,-0.3)
  {\textbf{POST}\\[2pt]\underline{id}\\ content\\ date\\ image\\ reply\_to\\ reply\_to\_parent};

% --- Associations ---
\node[assoc] (appart) at (2.6,0) {appartient};
\node[assoc] (publie) at (8.1,2.1) {publie};
\node[assoc] (aime)   at (8.1,-2.6) {aime\\{\scriptsize\underline{created\_at}}};
\node[assoc] (repond) at (12.6,3) {répond à};

% --- Liens + cardinalités ---
\draw[link] (role) -- (appart)
  node[card,pos=0.45,above]{0,n};
\draw[link] (appart) -- (user)
  node[card,pos=0.55,above]{1,1};

\draw[link] (user.north) |- (publie)
  node[card,pos=0.85,above]{0,n};
\draw[link] (publie) -| (post.north)
  node[card,pos=0.2,above]{1,1};

\draw[link] (user.south) |- (aime)
  node[card,pos=0.85,below]{0,n};
\draw[link] (aime) -| (post.south)
  node[card,pos=0.2,below]{0,n};

\draw[link] (post.north) to[out=110,in=20] (repond)
  node[card,pos=0.85,right]{0,n};
\draw[link] (repond) to[out=200,in=70] ([xshift=-6mm]post.north)
  node[card,pos=0.15,left]{0,1};
\end{tikzpicture}
}
\caption{Modèle conceptuel des données (formalisme Merise).}
\label{fig:mcd}
\end{figure}

\paragraph{Les réponses, sans table dédiée.} Une réponse n'est pas un type d'objet distinct :
c'est un \texttt{post} qui en référence un autre via \texttt{reply\_to} (le message auquel on
répond) et \texttt{reply\_to\_parent} (la racine du fil). Ce choix évite une table \texttt{responses}
séparée qui dupliquerait la structure d'un message, et permet de traiter messages et réponses avec
le même modèle et les mêmes requêtes.

\subsection{Modèle logique / physique (MLD, MPD)}

Le passage au modèle logique transforme les associations en clés étrangères. La règle de
nommage notable : le mot \texttt{user} étant \emph{réservé} en PostgreSQL, les colonnes de clé
étrangère se nomment \texttt{user\_id}. La figure~\ref{fig:mld} présente le schéma relationnel
et ses clés étrangères.

\begin{figure}[ht]
\centering
\resizebox{\textwidth}{!}{% MLD / schéma relationnel — tables physiques et clés étrangères
\begin{tikzpicture}[
  table/.style={draw=solageBlue!75, fill=infoBg, rounded corners=1pt,
    align=left, font=\ttfamily\scriptsize, inner sep=4pt},
  fk/.style={-{Stealth[length=2mm]}, solageGray, thin},
  fklab/.style={font=\scriptsize\ttfamily, text=todoOrange, fill=white, inner sep=1pt}
]
% --- Tables ---
\node[table] (roles) at (0,4)
  {\textbf{roles}\\ \underline{id}\\ name};
\node[table] (users) at (0,-0.4)
  {\textbf{users}\\ \underline{id}\\ name\\ firstname\\ email\\ password\\ \#\,role\\ image};
\node[table] (posts) at (6,1.6)
  {\textbf{posts}\\ \underline{id}\\ \#\,user\_id\\ date\\ content\\ reply\_to\\ reply\_to\_parent\\ image};
\node[table] (responses) at (12,2.2)
  {\textbf{responses}\\ \underline{id}\\ \#\,post\\ \#\,user\_id\\ date\\ content\\ \#\,reply\_to};
\node[table] (likes) at (6,-3.4)
  {\textbf{likes}\\ \underline{id}\\ \#\,user\_id\\ \#\,post\\ \#\,response\\ created\_at};
\node[table] (favs) at (0,-4.6)
  {\textbf{users\_favorites\_posts}\\ \underline{\#\,user\_id}\\ \underline{\#\,post}};

% --- Clés étrangères ---
\draw[fk] (users.north) -- (roles.south) node[fklab,midway]{role};
\draw[fk] (posts.west) -- (users.east) node[fklab,pos=0.5]{user\_id};
\draw[fk] (responses.south) to[out=-90,in=0] (posts.east) node[fklab,pos=0.6]{post};
\draw[fk] (responses) to[out=160,in=70] (users.north east) node[fklab,pos=0.85]{user\_id};
\draw[fk] (responses.north) to[out=120,in=60,looseness=6] (responses.east)
  node[fklab,pos=0.5,above]{reply\_to};
\draw[fk] (likes.west) to[out=180,in=-90] (users.south) node[fklab,pos=0.7]{user\_id};
\draw[fk] (likes.north) -- (posts.south) node[fklab,midway]{post};
\draw[fk] (likes.east) to[out=0,in=-90] (responses.south) node[fklab,pos=0.7]{response};
\draw[fk] (favs.north) to[out=90,in=180] (users.west) node[fklab,pos=0.6]{user\_id};
\draw[fk] (favs.east) to[out=0,in=-120] (posts.south west) node[fklab,pos=0.7]{post};
\end{tikzpicture}
}
\caption{Modèle logique des données : tables physiques et clés étrangères (\texttt{\#}).}
\label{fig:mld}
\end{figure}

Le modèle physique est concrétisé par le script de création \file{solage.pg.sql}, chargé une
seule fois à l'initialisation du conteneur PostgreSQL. L'extrait suivant en montre les tables
principales ; le script complet figure en annexe~\ref{ann:sql}.

\begin{lstlisting}[style=sql, caption={Extrait du script de création — \file{solage.pg.sql}}]
CREATE TABLE roles (
    id   SERIAL PRIMARY KEY,
    name VARCHAR(255) NOT NULL
);

CREATE TABLE users (
    id        SERIAL PRIMARY KEY,
    name      VARCHAR(255) NOT NULL,
    email     VARCHAR(255) NOT NULL UNIQUE,
    password  VARCHAR(255) NOT NULL,
    role      INT NULL REFERENCES roles(id) ON DELETE SET NULL,
    image     VARCHAR(255) NULL
);
CREATE INDEX users_role_idx ON users(role);

CREATE TABLE posts (
    id              SERIAL PRIMARY KEY,
    user_id         INT NULL REFERENCES users(id) ON DELETE CASCADE,
    date            TIMESTAMP NOT NULL,
    content         TEXT NULL,
    reply_to        INT NULL,
    reply_to_parent INT NULL,
    image           TEXT NULL
);
CREATE INDEX posts_user_idx ON posts(user_id);

CREATE TABLE likes (
    id         SERIAL PRIMARY KEY,
    user_id    INT NULL REFERENCES users(id) ON DELETE CASCADE,
    post       INT NULL REFERENCES posts(id) ON DELETE CASCADE,
    created_at TIMESTAMP NOT NULL
);
\end{lstlisting}

L'\textbf{intégrité} est garantie par les contraintes : \texttt{NOT NULL},
\texttt{UNIQUE(email)}, clés étrangères avec \texttt{ON DELETE CASCADE} (suppression en cascade
des messages d'un utilisateur supprimé) ou \texttt{SET NULL}. Des \textbf{index} accélèrent les
accès fréquents (clés étrangères, fil d'actualité).

\subsection{Évolution du schéma : migrations idempotentes}

Les évolutions additives ne sont pas appliquées à la main : elles passent par un système de
\textbf{migrations idempotentes} (\file{src/Migrations.php}), qui interroge
\texttt{information\_schema} avant d'agir et peut donc être rejoué sans risque.

\begin{lstlisting}[style=php, caption={Migration idempotente — \file{src/Migrations.php}}]
protected function columnExists(string $table, string $column): bool {
    $sql = "SELECT 1 FROM information_schema.columns
            WHERE table_schema = current_schema()
              AND table_name = :table AND column_name = :column LIMIT 1";
    $stmt = $this->db->prepare($sql);
    $stmt->execute([':table' => $table, ':column' => $column]);
    return $stmt->fetchColumn() !== false;
}

public function migrate(): void {
    $this->addColumnIfNotExists('posts', 'reply_to', 'INT', 'NULL');
    $this->addColumnIfNotExists('posts', 'reply_to_parent', 'INT', 'NULL');
    $this->addColumnIfNotExists('posts', 'image', 'TEXT', 'NULL');
}
\end{lstlisting}

\paragraph{Jeu d'essai et restauration.} Un jeu d'essai reproductible
(\file{seed.sql}) peuple une base de test (utilisateurs, messages, « j'aime », réponses).

\paragraph{Sauvegarde et restauration de la base.} La sauvegarde de la base de production est
automatisée par le script \file{scripts/backup.sh}, planifié en \emph{cron} (une fois par jour).
Il exécute \texttt{pg\_dump} \emph{à l'intérieur} du conteneur PostgreSQL — les identifiants ne
transitent jamais par la ligne de commande, ils sont lus dans l'environnement du conteneur — et
écrit un \emph{dump} compressé et horodaté dans \file{backups/}. L'option
\texttt{-{}-clean -{}-if-exists} rend la restauration rejouable (suppression puis recréation des
objets), un renommage atomique garantit qu'un \emph{dump} partiel n'est jamais conservé sous le
nom final, et les sauvegardes de plus de 14 jours sont purgées (rotation). Les \emph{dumps}
restent hors du dépôt Git (\file{backups/} est dans \file{.gitignore}) : ils contiennent des
données.

La restauration sur une base vierge — exigée par le critère C7 (« la base de test peut être
restaurée en cas d'incident ») — réinjecte le \emph{dump} décompressé dans \texttt{psql} :

\begin{lstlisting}[style=bash, caption={Restauration d'un \emph{dump} compressé dans le conteneur}]
gunzip -c backups/solage-AAAA-MM-JJ-HHMMSS.sql.gz | \
  docker compose -f docker-compose.prod.yml exec -T postgres \
  sh -c 'psql -U "$POSTGRES_USER" -d "$POSTGRES_DB"'
\end{lstlisting}

\section{Diagrammes UML}
\label{sec:uml}

\subsection{Diagramme de cas d'utilisation}

La figure~\ref{fig:usecase} formalise le comportement attendu : les trois acteurs et leurs cas
d'usage, avec la généralisation Administrateur $\rightarrow$ Utilisateur $\rightarrow$ Visiteur.

\begin{figure}[ht]
\centering
\resizebox{0.92\textwidth}{!}{% Diagramme de cas d'utilisation
\begin{tikzpicture}[
  uc/.style={draw=solageBlue!70, fill=infoBg, ellipse, align=center,
    font=\scriptsize, minimum width=3cm, minimum height=0.85cm, inner sep=1pt},
  lnk/.style={solageGray, thin},
  gen/.style={-{Triangle[open,length=3mm,width=3mm]}, solageDark, semithick},
  actorlbl/.style={font=\scriptsize\bfseries, align=center},
  actor/.pic={
    \draw[thick] (0,0) circle (3mm);
    \draw[thick] (0,-3mm) -- (0,-1.2);
    \draw[thick] (-5mm,-0.75) -- (5mm,-0.75);
    \draw[thick] (0,-1.2) -- (-5mm,-1.8);
    \draw[thick] (0,-1.2) -- (5mm,-1.8);
  }
]
% --- Acteurs ---
\pic at (0,2) {actor};      \node[actorlbl] at (0,-0.2) {Visiteur};
\pic at (0,-2.6) {actor};   \node[actorlbl] at (0,-4.8) {Utilisateur};
\pic at (0,-7.2) {actor};   \node[actorlbl] at (0,-9.4) {Administrateur};

% --- Frontière du système ---
\node[draw=solageDark, rounded corners, minimum width=7cm, minimum height=12.4cm,
      label={[font=\bfseries]north:Solage}] (sys) at (7.3,-3.2) {};

% --- Cas d'utilisation ---
\node[uc] (inscr)  at (7.3,2.1)  {S'inscrire};
\node[uc] (conn)   at (7.3,1.1)  {Se connecter};
\node[uc] (fil)    at (7.3,0.1)  {Consulter le fil};
\node[uc] (pub)    at (7.3,-0.9) {Publier un message};
\node[uc] (rep)    at (7.3,-1.9) {Répondre à un message};
\node[uc] (aime)   at (7.3,-2.9) {Aimer un message};
\node[uc] (rech)   at (7.3,-3.9) {Rechercher};
\node[uc] (prof)   at (7.3,-4.9) {Éditer son profil};
\node[uc] (suppr)  at (7.3,-5.9) {Supprimer son contenu};
\node[uc] (moder)  at (7.3,-8.6) {Modérer (tout contenu)};

% --- Associations acteur / cas ---
\draw[lnk] (0,1.3) -- (inscr.west);
\draw[lnk] (0,1.3) -- (conn.west);
\foreach \c in {fil,pub,rep,aime,rech,prof,suppr}
  \draw[lnk] (0,-3.3) -- (\c.west);
\draw[lnk] (0,-7.9) -- (moder.west);

% --- Généralisations entre acteurs (décalées à gauche pour ne pas recouvrir) ---
\draw[gen] (-1.2,-4.4) -- (-1.2,0.0);          % Utilisateur -> Visiteur
\draw[gen] (-1.2,-9.0) -- (-1.2,-4.8);         % Admin -> Utilisateur
\node[font=\tiny,text=solageGray,rotate=90] at (-1.5,-2.2) {généralise};
\node[font=\tiny,text=solageGray,rotate=90] at (-1.5,-6.9) {généralise};
\end{tikzpicture}
}
\caption{Diagramme de cas d'utilisation de \projet{}.}
\label{fig:usecase}
\end{figure}

\subsection{Diagramme de séquence : « Publier un message »}

Le cas le plus significatif est la publication d'un message : il traverse toutes les couches et
mobilise les protections de sécurité. La figure~\ref{fig:sequence} en détaille le déroulé, de
l'appel \texttt{fetch} du navigateur jusqu'à l'insertion en base et au rendu optimiste.

\begin{figure}[ht]
\centering
\resizebox{\textwidth}{!}{% Diagramme de séquence — « Publier un message »
\begin{tikzpicture}[
  font=\scriptsize,
  ll/.style={draw=solageBlue!70, fill=infoBg, rounded corners=1pt,
    minimum width=1.9cm, minimum height=0.7cm, align=center, font=\scriptsize\bfseries},
  msg/.style={-{Stealth[length=2mm]}, semithick},
  ret/.style={-{Stealth[length=2mm]}, semithick, dashed, solageGray},
  lbl/.style={font=\scriptsize, fill=white, inner sep=1.5pt},
  slbl/.style={font=\scriptsize, anchor=west, inner sep=1.5pt}
]
% --- Lignes de vie ---
\node[ll] (nav)  at (0,0)    {Navigateur\\[-1pt]{\scriptsize index.js}};
\node[ll] (rt)   at (2.9,0)  {Router};
\node[ll] (mw)   at (5.8,0)  {Middlewares\\[-1pt]{\scriptsize CSRF · Auth}};
\node[ll] (ctrl) at (8.9,0)  {PostController};
\node[ll] (mdl)  at (11.7,0) {PostModel};
\node[ll] (db)   at (14.2,0) {PostgreSQL};
\foreach \n in {nav,rt,mw,ctrl,mdl,db} \draw[dashed,solageGray!55] (\n.south) -- ++(0,-11.4);

% --- Messages (yshift en mm depuis l'ancre .south des en-têtes) ---
\draw[msg] ([yshift=-9mm]nav.south) -- ([yshift=-9mm]rt.south)
  node[lbl,midway,above]{POST /api/post {\scriptsize(FormData + X-CSRF-Token)}};
\draw[msg] ([yshift=-18mm]rt.south) -- ([yshift=-18mm]mw.south)
  node[lbl,midway,above]{handle() : jeton CSRF puis session};
\draw[ret] ([yshift=-26mm]mw.south) -- ([yshift=-26mm]rt.south)
  node[lbl,midway,above]{OK {\scriptsize(sinon 403 / redirection)}};
\draw[msg] ([yshift=-35mm]rt.south) -- ([yshift=-35mm]ctrl.south)
  node[lbl,midway,above]{create()};
\draw[msg] ([yshift=-43mm]ctrl.south) to[out=-35,in=35] ([yshift=-50mm]ctrl.south);
\node[slbl] at ([xshift=6mm,yshift=-46mm]ctrl.south) {valider upload + contenu};
\draw[msg] ([yshift=-58mm]ctrl.south) -- ([yshift=-58mm]mdl.south)
  node[lbl,midway,above]{createPost()};
\draw[msg] ([yshift=-67mm]mdl.south) -- ([yshift=-67mm]db.south)
  node[lbl,midway,above]{INSERT {\scriptsize(requête préparée)}};
\draw[ret] ([yshift=-75mm]db.south) -- ([yshift=-75mm]mdl.south)
  node[lbl,midway,above]{id du post};
\draw[ret] ([yshift=-83mm]mdl.south) -- ([yshift=-83mm]ctrl.south)
  node[lbl,midway,above]{id};
\draw[ret] ([yshift=-91mm]ctrl.south) -- ([yshift=-91mm]nav.south)
  node[lbl,midway,above]{JSON \{success, data\}};
\draw[msg] ([yshift=-99mm]nav.south) to[out=-35,in=35] ([yshift=-106mm]nav.south);
\node[slbl] at ([xshift=6mm,yshift=-102mm]nav.south) {rendu optimiste {\scriptsize(escapeHtml)}};
\end{tikzpicture}
}
\caption{Diagramme de séquence : publication d'un message (avec passage par les middlewares
CSRF et authentification).}
\label{fig:sequence}
\end{figure}

\subsection{Diagramme de classes (couche modèle)}

La figure~\ref{fig:classes} présente un extrait de la couche modèle et le gestionnaire de
session, avec leurs principales méthodes et dépendances.

\begin{figure}[ht]
\centering
\resizebox{\textwidth}{!}{% Diagramme de classes (extrait : couche modèle + session)
\begin{tikzpicture}[
  class/.style={draw=solageBlue!75, fill=infoBg, rectangle split,
    rectangle split parts=3, align=left, font=\scriptsize, inner sep=3pt,
    text width=4.3cm},
  assoc/.style={-{Stealth[length=2mm]}, semithick, solageGray},
  dep/.style={-{Stealth[length=2mm]}, semithick, dashed, solageGray},
  albl/.style={font=\scriptsize\itshape, text=solageDark, fill=white, inner sep=1pt}
]
\node[class] (user) at (0,0) {
  \textbf{UserModel}
  \nodepart{second}
  - id, name, email\\ - password, role, image
  \nodepart{third}
  + getUserById(id)\\ + getUserByEmail(email)\\ + getUsersByIds(ids[]): map\\ + createUser(...)\\ + getRoleName(): string
};
\node[class] (post) at (-5.6,-0.2) {
  \textbf{PostModel}
  \nodepart{second}
  - id, user, content\\ - date, image, replyTo
  \nodepart{third}
  + getPosts(): array\\ + getPostById(id)\\ + createPost(...)\\ + getResponses(): array\\ + delete(id): bool
};
\node[class] (like) at (-5.6,-5.4) {
  \textbf{LikeModel}
  \nodepart{second}
  - id, userId\\ - post, response
  \nodepart{third}
  + create(...)\\ + isLikedBy(...): bool\\ + countForPost(id): int
};
\node[class] (sess) at (5.6,-0.6) {
  \textbf{SessionManager}
  \nodepart{second}
  - user: UserModel
  \nodepart{third}
  + login(id)\\ + isLoggedIn(): bool\\ + isAdmin(): bool\\ + getUserId(): int
};

\draw[assoc] (sess) -- (user) node[albl,midway,above]{utilise};
\draw[dep] (post.north west) to[out=120,in=60] (user.north)
  node[albl,pos=0.5,above]{user\_id};
\draw[dep] (like.east) -- (post.south east|-like.east)
  node[albl,midway,below]{post};
\draw[dep] (like.north) -- (post.south) node[albl,midway,right]{};
\end{tikzpicture}
}
\caption{Diagramme de classes (extrait) : modèles et gestion de session.}
\label{fig:classes}
\end{figure}

